\chapter{Software Construction}
\label{chap:construction}

Software construction, as defined in Lecture~C08, is the activity of converting a
detailed design into working, maintainable code. It encompasses not only typing
instructions into an editor but also the discipline that surrounds that act: a
consistent coding style, a modular decomposition that hides complexity behind stable
interfaces, a principled programme of refactoring to remove code smells as they
accumulate, and a formal inspection that exposes defects before they reach users. This
chapter traces each of those concerns through the SHEO implementation, grounding every
point in the structure of the system as it was actually built.

\section{Construction in the development lifecycle}
\label{sec:construction-lifecycle}

In the SHEO development process, construction follows the detailed design established
in Chapter~\ref{chap:design}: class responsibilities, sequence flows, and persistence
schemas had been decided before implementation began. Construction therefore had a
clear contract to satisfy, and every significant decision---which library, which
language feature, how to structure an error return---could be evaluated against the
design goals of high cohesion, low coupling, and information hiding that are
documented in the System Design Document.

The technology stack was chosen to support those goals directly.

\begin{table}[htbp]
\centering
\caption{SHEO technology stack and the construction rationale for each choice.}
\label{tab:stack}
\begin{tabularx}{\textwidth}{L{0.22\textwidth} L{0.18\textwidth} Y}
\toprule
\textbf{Component} & \textbf{Technology} & \textbf{Construction rationale} \\
\midrule
Language & A modern, statically hinted dynamic language &
  Its rich type-hint system, value-object facilities, and a mature standard library
  reduce the volume of boilerplate that obscures intent. \\
Web framework & A declarative Python web framework &
  Routes, dependency injection, and request and response schemas are declared
  together, making the interface layer self-documenting without a separate tool. \\
Persistence & A mature object-relational mapper &
  The repository pattern is natural to it; queries are isolated from
  service logic, realising the information-hiding goal at the persistence boundary. \\
Data validation & A type-driven validation library &
  Schema validation is declared as ordinary types; invalid inputs are rejected at the
  boundary with a machine-readable error before any business logic executes. \\
Relational store & A file-backed embedded database &
  Eliminating an external database server removes an infrastructure dependency during
  development and testing, keeping the test suite fast and reproducible. \\
\midrule
User-interface framework & A component-based front-end framework &
  Component-based decomposition mirrors the backend service decomposition, giving
  each screen a single owner and a single reason to change. \\
Charting & A declarative charting library &
  Declarative, data-driven charts align with the capability-schema philosophy: the
  dashboard renders from data descriptors rather than hard-coded widget trees. \\
\bottomrule
\end{tabularx}
\end{table}

\section{Coding style and conventions}
\label{sec:coding-style}

Lecture~C08 identifies six coding-style principles: \emph{consistency},
\emph{clarity}, \emph{single-purpose methods}, \emph{fewer than six parameters},
\emph{naming} (verbs for functions, nouns for classes and variables), and
\emph{comments that describe what, not how}. SHEO adheres to each, enforced by
convention rather than a separate lint pass.

\subsection{Naming and intent}

Every operation in the backend is named by a verb phrase that announces its effect:
the capability-validation routine, the command-application entry point, the
unit-scoped lookup that raises a not-found error when nothing matches, the
capability-retrieval accessor, and the adapter-selection function all read as actions.
Every abstraction is named by a noun that describes the responsibility it represents:
the device-control service, the capability schema, the device-adapter interface, and
the structured validation result. This consistency means that a reader surveying the
list of components in a module can determine the responsibility of each one before
examining its definition.

\subsection{File-level docstrings}

Every backend module opens with a brief descriptive header that states the module's
purpose and its relationship to the requirements. The header of the device-control
service module, for example, records that the module implements the capability-driven
control path of \req{REQ-4.2.3} and the safety guards of \req{NFR-SAF}, and that both
manual user commands and rule-initiated commands flow through the single
command-application entry point so that validation and safety rules are enforced
uniformly. Such a header names the requirement identifiers that the module satisfies
and states the architectural invariant---the single choke point for all commands---that
reviewers should verify, in accordance with the ``describe what, not how'' rule.

\subsection{Type hints and small, focused functions}

Type annotations are declared on every public operation. Return types,
parameter types, and the element types of collections are all stated, converting the
type system into a machine-checked expression of each operation's contract. The
capability-validation routine is representative of the single-purpose, narrow-parameter
style the course advocates.

\label{lst:validate}%
The routine accepts a device type, the name of a control, and a proposed value, and
returns a structured validation result that reports either success---with the value
normalised to its canonical form---or failure together with a human-readable
explanation. It first resolves the capability schema for the device type, then locates
the named control within that schema, and finally checks the proposed value against the
control's kind: a numeric range control is bounds-checked, while toggle and enumeration
controls are matched against their permitted sets. The routine takes exactly three
parameters, well below the six-parameter threshold; it returns a structured result so
that the error message is machine-readable rather than raised as a bare exception; and
it has no side effects, leaving persistence and dispatch to its callers. Its descriptive
header cites \req{REQ-4.2.3} rather than restating the algorithm that the code already
expresses.

\section{Information hiding and modularity in code}
\label{sec:info-hiding}

Information hiding is applied at three distinct boundaries in the SHEO codebase, each
chosen to isolate a design decision that is likely to change.

\subsection{Repositories hide persistence}

The persistence layer is organised into repositories---one for devices and one for the
per-device permission records---that encapsulate all query construction. Service code
never assembles a query predicate directly; it invokes named operations such as a
unit-scoped query that returns every device belonging to a given unit, or a membership
check that confirms a particular device belongs to a particular unit, and receives
domain objects in return. If the storage engine were replaced, only the repository
implementations would need to change.

\subsection{The capability schema hides device-type variance}

The capability schema expresses everything the rest of the system needs to know about a
device type as \emph{data} rather than as \emph{code branches}. Neither the interface
layer nor the user interface ever tests for a specific device type with a conditional
branch; both query the schema and act on what they find. Adding a new device type
therefore requires changes in exactly two places---the capability registry and the
corresponding adapter---and nowhere else.

\subsection{Adapters hide per-type behaviour}

The Strategy pattern is realised through an abstract device-adapter interface defined
in the adapter base module. The interface is deliberately minimal: one abstract
operation for advancing the simulation by one time step, and one default operation for
applying a validated control change to a device's state. Callers obtain the correct
strategy through the adapter-selection function, whose body is a single lookup in a
registry that maps each device type to its corresponding per-type adapter; the concrete
adapter classes are never referenced directly by callers.

\label{lst:adapter}%
The consequence of this design is that the entire call site in the service layer
reduces to two steps---obtain the adapter for the device's type, then ask it to apply
the validated control change---regardless of how complex the underlying simulation logic
becomes. The air-conditioner adapter, for instance, models hysteresis, compressor
cycling, and ambient-temperature drift in roughly thirty lines that remain entirely
invisible to the service.

\section{Refactoring and code smells}
\label{sec:refactoring}

Refactoring, as defined in C08, is the transformation of the internal structure of
code without changing its observable behaviour. Its purpose is to remove \emph{code
smells}---structural problems that make software harder to understand, test, or
extend---prior to the addition of new capability. The canonical smells named in the
course include duplicated code, long method, large class, too many parameters,
divergent change, lazy class, speculative generality, and temporary field.

Two concrete refactoring decisions made during SHEO construction directly address
smells from that list.

\subsection{Centralising the safety check}

An early version of the control path performed safety checks in multiple places: once
in the interface handler that received a user command, and again in the rule engine
that executed automated actions. This is the \emph{duplicated code} smell:
the same guard appeared in two locations, so a correction to the guard had to be
applied twice and could silently diverge. The refactoring moved all safety logic into
a single private safety-check step within the device-control service, and introduced a
single public command-application entry point through which every command---manual and
automated alike---must pass.

\label{lst:apply}%
Within that entry point the command-application logic is strictly linear. It first
rejects commands addressed to a device that is offline or has been forced offline; it
then delegates to the capability-validation routine and abandons the command if the
proposed value is invalid, retaining the normalised value otherwise; it next runs the
safety-check step and returns immediately if a safety constraint is violated; and only
then does it obtain the per-type adapter and ask it to apply the validated change. The
sequence---validate, safety-check, dispatch---is explicit and self-contained. A reader
can verify the entire guard chain by following one operation, without tracing branches
across several modules.

\subsection{Removing duplicated validation}

Before the capability schema was introduced, each interface operation that accepted a
device command contained its own range check and enumeration check. These checks were
written independently, carried slightly different error messages, and diverged over time
as device types were added. This is simultaneously the \emph{duplicated code} smell and
a mild case of the \emph{divergent change} smell: a new device type required edits
across several interface operations. The refactoring introduced the capability-validation
routine as the single authority on what constitutes a valid command for each device
type, and removed all per-operation checks. Every operation now delegates to that
routine by way of the command-application entry point; the capability schema is the
single source of truth.

Both refactorings satisfy the C08 criterion: no behaviour changed---the automated test
suite continued to pass at every intermediate step---but the internal structure became
easier to reason about and to extend.

\section{Code review}
\label{sec:code-review}

Lecture~C08 distinguishes two review processes. A \textbf{walkthrough} is author-led
and informal: the author guides peers through the code, gathering suggestions. A
\textbf{formal inspection} is conducted by non-authors against a pre-defined
checklist; the author is present but does not lead. The course identifies three review
criteria: \emph{maintainability} (can the code be changed safely?),
\emph{reusability} (can components be used in other contexts?), and
\emph{extensibility} (can the system accommodate new requirements without pervasive
edits?).

\subsection{The adversarial inspection process}

SHEO underwent a structured review that functioned as a formal inspection. A
non-author reviewing agent examined the codebase against a checklist derived from the
requirements, the non-functional requirements, and security best practice. The review
was adversarial in the sense that the inspector was explicitly tasked with finding
defects, not confirming correctness---matching the spirit of the formal inspection
model in which the non-author role is designed to surface problems that the author has
normalised. The checklist covered five categories:

\begin{enumerate}
  \item \textbf{Authorisation.} Every interface operation must enforce unit-scoped
        access control; cross-unit access must be impossible.
  \item \textbf{Input validation.} Every external input must be rejected at the
        boundary before entering business logic; no operation may bypass the
        validation schema.
  \item \textbf{Safety constraint enforcement.} The \req{NFR-SAF} rules must be
        applied on every command path---interactive and automated---not only when a
        human issues a command.
  \item \textbf{Telemetry isolation.} One unit must not be able to read another
        unit's sensor data, even if device identifiers are predictable integers.
  \item \textbf{Rule mutation.} A rule update must pass through the same
        authorisation and validation as rule creation.
\end{enumerate}

\subsection{Defects found and corrected}

The inspection identified two defects, both verified by a failing test before the fix
and a passing test after.

\paragraph{Cross-unit insecure direct object reference in the telemetry series
operation.}
The telemetry-series retrieval operation accepted a device identifier without verifying
that the device belonged to the requesting user's unit. An authenticated resident from
unit~A could therefore obtain the full telemetry history of any device in unit~B by
supplying its identifier. This is an insecure direct object reference (IDOR)
vulnerability---a failure of the authorisation criterion. The fix introduced a
unit-scoped lookup at the start of the operation, matching the pattern already
established by the device-control service's unit-scoped lookup. A regression test was
added asserting that a not-found response is returned when a device belonging to a
different unit is requested.

\paragraph{Unvalidated rule update.}
The rule-update operation accepted an arbitrary request payload and wrote it to the
database without passing it through the same input-validation schema that the
rule-creation operation enforced. A malformed or adversarially crafted update could
therefore insert a rule with missing required fields, causing the rule engine to raise
an unhandled exception at execution time. The fix applied the existing rule-update
schema to the operation, restoring the input-boundary guarantee established at creation
time.

\subsection{Review criteria mapped to SHEO mechanisms}

Table~\ref{tab:review-criteria} maps each C08 review criterion to the mechanism
that satisfies it after the inspection.

\begin{table}[htbp]
\centering
\caption{C08 code review criteria and the corresponding SHEO mechanisms
post-inspection.}
\label{tab:review-criteria}
\begin{tabularx}{\textwidth}{L{0.22\textwidth} Y Y}
\toprule
\textbf{Review criterion} & \textbf{What it requires} & \textbf{SHEO mechanism} \\
\midrule
Maintainability &
  Internal structure can be changed without inadvertent side effects. &
  A single command-application choke point; the repository abstraction; one
  capability schema per device type; no duplicated guards. \\
Reusability &
  Components can be used in different contexts without modification. &
  The capability-validation routine is invoked by the interface layer, the rule
  engine, and the automated test suite without any context dependency; the per-type
  adapters are stateless and context-free. \\
Extensibility &
  New requirements can be accommodated with localised changes. &
  A new device type requires exactly one new capability-schema entry and one new
  per-type adapter; no existing code is modified (Open/Closed Principle). \\
Security &
  Access control is enforced at every boundary. &
  A unit-scoped lookup on every device and telemetry operation; input-validation
  on every write path; an administrator-authorisation guard on all mutation
  operations. \\
\bottomrule
\end{tabularx}
\end{table}

The two defects found by the inspection were both failures of the
\emph{maintainability} and \emph{security} criteria: the cross-unit IDOR arose
because a defensive pattern established in one service had not been propagated to a
newer interface operation, and the unvalidated rule update arose because a boundary
guarantee had been applied inconsistently. Both are the kind of defect that is invisible to the
original author---who knows the intent---and visible to a non-author reviewer who
tests the code against the stated contract. This is precisely the rationale that C08
gives for requiring formal inspections rather than relying on walkthroughs alone.
